%---------------------------------------------------------------
% Kurzfassung
%---------------------------------------------------------------
\begin{kurzfassung}
Die Ultraschalltechnologie eignet sich in zahlreichen Prozessen und Applikationen schon lange ideal für den automatisierten Einsatz. Sie bietet Lösungen sowohl für einfache als auch für anspruchsvolle Schweiß- und Trennaufgaben an Einzelarbeitsplätzen sowie in Produktionslinien und in Sonderanlagen. Typische Einsatzgebiete sind das Schweißen von Kunststoffen und Buntmetallen, das Trennschweißen von Vliesen, Portionieren von Lebensmitteln und viele weitere Anwendungen. Das Herzstück in dieser Anwendungen ist der Ultraschallkonverter, der die elektrische Wechselspannung in eine nötige Schwingungsamplitude umwandelt.  Für die Bearbeitung ist diese Schwingungsamplitude maßgebend. Da diese Schwingungsamplitude einen Einfluss auf die gesamte Prozessqualität, die Prozessgeschwindigkeit und die Lebensdauer des Ultraschallkonverters hat.
Im Rahmen dieser Arbeit wird ein neuartiger Ultraschallkonverter behandelt, der mit zwei integrierten Sensorelementen ausgestattet ist. Diese integrierten Sensorelementen sind aus zwei unterschiedlichen Materialen hergestellt, die jeweils unterschiedliche elektrische sowie mechanische Eigenschaften aufweisen.  Um den Bearbeitungsprozess mit konstanter Qualität ausführen zu können, soll die Schwingungsamplitude gemessen und durch Anpassen der Amplitude der angelegten Wechselspannung geregelt werden.
Dazu ist jedoch eine Signalaufbereitung notwendig, die die verschiedenen „roh“-Signale der Sensorscheiben miteinander kombiniert, um die Informationen wie Amplitude und Temperatur zu extrahieren. Diese Signalaufbereitung erfolgt durch Filterung der roh-Signale, die in dieser Arbeit ausführlich erörtert wird. 
Eine besondere Herausforderung besteht darin, die Signale möglichst rauscharm von den Sensorscheiben zum einige Meter entfernten Messgerät zu transportieren. Die Amplitude des Messignals der Sensorkeramiken ist mit einigen Volt sehr gering im Vergleich zu den überlagerten Spannungen von über 1kV, die durch den piezoelektrischen Effekt aufgrund der mechanischen Schwingungen überlagert sind. Zur bestmöglichen Lösung dieses Problems werden verschiedene Methoden der passiven und aktiven Filterung evaluiert, aufgebaut und getestet werden. Die gefilterten Signale werden anschließend möglichst rauscharm vom Konverter zum einige Meter entfernten Messgerät übertragen. Auch hierfür werden verschiedene Signalübertragungsmethoden evaluiert.
\vspace{1cm}


{\bfseries Schlagwörter:} bla, bla , bla, bla, bla

\end{kurzfassung}
